\phantomsection
\color{unidarkgreen}\section*{\centering{\large{ПЕРЕЧЕНЬ СОКРАЩЕНИЙ}}}
\addcontentsline{toc}{section}{Перечень сокращений}
\noindent\makebox[0pt][l]{\smash{\color{white}\fontsize{0.01pt}{0.01pt}\selectfont<begABBRS>}}%
\color{black}
\begin{longtable}{>{\raggedright\arraybackslash}m{2cm}>{\raggedright\arraybackslash}m{0.5cm}>{\raggedright\arraybackslash}m{20cm}}
\endfirsthead\endhead\endfoot\endlastfoot
IEC & -- &  International Electrotechnical Commission; \\
PPS & -- & Pulse Per Second; \\
АВР & -- & автоматическое включение резерва; \\
АПВ & -- & автоматическое повторное включение; \\
БИ & -- & блок испытательный; \\
БНН & -- & блокировка при неисправности цепей напряжения; \\
ВН & -- & высшее напряжение; \\
ГЗ & -- & газовая защита; \\
ГЗТ & -- & газовая защита трансформатора; \\
ДВ & -- & дискретный вход; \\
ЕЭС & -- & единая энергетическая система; \\
ЗНР & -- & защита от неполнофазного режима; \\
ЗНФ & -- & защита от непереключения фаз; \\
ИО & -- & измерительный орган / избирательный орган; \\
ИЧМ & -- & интерфейс человек-машина; \\
ИЭУ & -- & интеллектуальное электронное устройство; \\
КИ & -- & контроль изоляции; \\
НКУ & -- & низковольтное комплектное устройство; \\
НН & -- & низшее напряжение; \\
НПП & -- & научно-производственное предприятие; \\
НТД & -- & нормативно-техническая документация; \\
ОВ & -- & оперативный вывод / обходной выключатель; \\
ОПУ & -- & общеподстанционный пункт управления; \\
ОТ & -- & оперативный ток; \\
ПК & -- & предохранительный клапан; \\
ПО & -- & пусковой орган / программное обеспечение; \\
ПС & -- & предупредительная сигнализация; \\
ПУЭ & -- & правила устройства электроустановок; \\
РАС & -- & регистратор аварийных событий; \\
РЗА & -- & релейная защита и автоматика; \\
РЗТ & -- & резервные защиты трансформатора; \\
РПВ & -- & реле положения «включено»; \\
РПН & -- & устройство регулирования напряжения под нагрузкой; \\
РПО & -- & реле положения отключено; \\
РФ & -- & реле фиксации; \\
РЭ & -- & руководство по эксплуатации; \\
СВ & -- & секционный выключатель; \\
СИ & -- & средства измерений; \\
СН & -- & среднее напряжение; \\
СШ & -- & система (секция) шин; \\
ТЗНП & -- & токовая защита нулевой последовательности; \\
ТК & -- & токовый контроль / текущий контроль; \\
ТН & -- & трансформатор напряжения; \\
ТО & -- & токовая отсечка; \\
ТУ & -- & телеуправление / телеускорение / технические условия; \\
УРОВ & -- & устройство резервирования при отказе выключателя; \\
ЦН & -- & цепи напряжения; \\
ШСВ & -- & шиносоединительный выключатель; \\
ШЭТ & -- & шкаф электротехнический типовой; \\
ЭМО & -- & электромагнит отключения. \\
{\color{white}\fontsize{0.1pt}{0.1pt}\selectfont<endABBRS>}
\end{longtable}